\section{Nullspace Bias Perturbation}
\label{sec:method}
\subsection{Setting and notation}
Let $\theta_0$ be a public pretrained checkpoint and $\theta_e$ a fine-tuned expert. The defender publishes a modified checkpoint $\theta_e'$ and seeks to retain standalone utility while reducing protected-task utility after parameter merging. Construction uses only $\theta_0$, $\theta_e$, and a fixed task classifier. The defender does not inspect a partner checkpoint or task images.

The downstream user has the released checkpoint and the public base. We test direct merging and published recovery procedures, and separately examine a method-specific readout cleanup. Our objective is to disrupt unauthorized merging; it does not require preserving companion-task accuracy in the merged model. The released expert must retain its standalone function.

Write the final visual embedding of the expert as
\begin{equation}
 t_e(x)=[z_e(x)\operatorname{diag}(\gamma_e)+b_e^{\mathsf T}]P_e,
 \label{eq:embedding}
\end{equation}
Here $P_e\in\R^{d\times p}$ is the visual projection. Hidden features $z_e\in\R^{1\times d}$ and embeddings $t_e\in\R^{1\times p}$ are row vectors; $b_e,\gamma_e\in\R^d$ are column vectors. The features $z_e$ are centered and variance-normalized before the learned LayerNorm affine transformation. For the evaluated ViT-B/32 and ViT-B/16 encoders, $d=768$, $p=512$, and $P_e$ has full column rank. With the usual positive LayerNorm epsilon,
\begin{equation}
 z_e\mathbf{1}_d=0,\qquad \|z_e\|_2\leq\sqrt d.
 \label{eq:lnbound}
\end{equation}
The classifier uses fixed class vectors $w_j\in\R^p$, zero class biases and a shared positive scale $\eta$:
\begin{equation}
 s_j(x)=\eta\,\frac{t_e(x)w_j}{\|t_e(x)\|_2}.
 \label{eq:cosine}
\end{equation}
For nonzero embeddings, class comparisons therefore depend only on the signs of $t_e(x)(w_k-w_j)$.

\begin{figure}[!t]
\centering
\begin{tikzpicture}[font=\footnotesize,>=Latex,
 node distance=5mm,
 box/.style={draw,rounded corners=1pt,align=center,minimum height=8mm,inner sep=3pt},
 arr/.style={->,line width=.55pt}]
\node[anchor=west] at (-.15,2.8) {\textbf{Standalone expert}};
\node[box] (ln) at (.65,2.15) {Final LN\\$b_e+\tau u$};
\node[box,right=9mm of ln] (proj) {$P_e$};
\node[box,right=9mm of proj] (out) {Same\\embedding};
\draw[arr] (ln)--(proj);
\draw[arr] (proj)--(out);
\node[below=2mm of proj,align=center] {$u^{\mathsf T}P_e=0$};
\node[anchor=west] at (-.15,.8) {\textbf{Merged model}};
\node[box] (mln) at (.65,.15) {Merged LN\\$b_m+\lambda_e\tau u$};
\node[box,right=9mm of mln] (mproj) {$P_m$};
\node[box,right=9mm of mproj] (mout) {Changed\\embedding};
\draw[arr] (mln)--(mproj);
\draw[arr] (mproj)--(mout);
\node[below=2mm of mproj,align=center] {$\Delta t_m=\lambda_e\tau u^{\mathsf T}P_m$};
\end{tikzpicture}
\caption{The bias perturbation vanishes under the expert's own projection. Merging changes the projection and can expose a constant feature offset. The actual classification effect is measured for each fusion rule and partner. This is a structural illustration.}
\label{fig:mechanism}
\end{figure}

\subsection{Preserving the standalone encoder}
Let $N_e\in\R^{d\times(d-p)}$ be an orthonormal basis of $\ker(P_e^{\mathsf T})$. As shown in \cref{fig:mechanism}, choose $u=N_ea$ and publish
\begin{equation}
 b_e'=b_e+\tau u,\qquad \tau>0,
 \label{eq:bias}
\end{equation}
leaving all other parameters unchanged.

\begin{proposition}[Standalone preservation]
If $P_e^{\mathsf T}u=0$, then the modification in \cref{eq:bias} preserves the visual embedding in exact arithmetic for every input.
\end{proposition}
\begin{proof}
The LayerNorm input, mean, variance, gain, and visual projection are unchanged. Hence the embedding difference is $\tau u^{\mathsf T}P_e=0$.
\end{proof}

The identity applies to every input and any fixed head on the visual embedding. A nonzero construction requires a nontrivial left nullspace; a full-row-rank projection with $d\le p$ does not provide one. We use an SVD to compute the basis and store the released bias in FP32. This leaves the architecture and inference cost unchanged. Large bias values can increase cancellation error in floating point, which we assess on the full standalone test sets.

\subsection{Why merging can expose the perturbation}
For Task Arithmetic with coefficients $\lambda_i$, the merged checkpoint is
\begin{equation}
 \theta_m=\theta_0+\sum_i\lambda_i(\theta_i-\theta_0).
 \label{eq:ta}
\end{equation}
If expert $e$ is protected, its bias perturbation enters with coefficient $\lambda_e$. Denote the merged visual projection by $P_m$. Relative to the same merged model without this bias update, the embedding difference is exactly
\begin{equation}
 \Delta t_m(x)=\lambda_e\tau u^{\mathsf T}P_m.
 \label{eq:exposure}
\end{equation}
The offset is constant across inputs because it is added after hidden-feature normalization. To separate the available public information from the unknown partners, write $\Delta P_i=P_i-P_0$. Using $u^{\mathsf T}P_e=0$ gives
\begin{equation}
\begin{aligned}
 u^{\mathsf T}P_m
 &=u^{\mathsf T}\widetilde P(\lambda_e)+u^{\mathsf T}E,\\
 \widetilde P(\lambda_e)&=(1-\lambda_e)P_0+\lambda_eP_e,\\
 E&=\sum_{i\ne e}\lambda_i\Delta P_i.
\end{aligned}
\label{eq:partners}
\end{equation}
The first term is available before release. The second describes the contribution of the partners' task updates. This exact decomposition motivates using the public-base interpolation to choose the perturbation. It does not assume that $E$ is small. For a class contrast $v=w_k-w_j$, the discrepancy satisfies
\begin{equation}
 |u^{\mathsf T}(P_m-\widetilde P)v|
 \leq \|u\|_2\|Ev\|_2.
 \label{eq:partnererror}
\end{equation}
Thus a positive proxy offset remains positive whenever its gap exceeds this error, although a constant prediction also requires the full feature bound below. For two experts at coefficient $0.5$, the offset reduces to $\tau u^{\mathsf T}P_{\mathrm{partner}}/4$. A partner with the same left nullspace exposes no perturbation. We evaluate partner models rather than assuming a universal transfer guarantee.

\subsection{Choosing a direction without images}
We use a public-base interpolation with fixed coefficient $\lambda=0.5$ to select a direction. This interpolation uses the protected expert and the public base, without a partner expert. Denote its final parameters by $\widetilde P,\widetilde\gamma,\widetilde b$. Let $W$ contain the fixed class vectors as rows and define the centered classifier $W_c=(I_C-\mathbf{1}_C\mathbf{1}_C^{\mathsf T}/C)W$.

Compute the thin SVD
\begin{equation}
 H=N_e^{\mathsf T}\widetilde P W_c^{\mathsf T}=U\Sigma V^{\mathsf T}.
\end{equation}
A fixed seed selects a class index before evaluation. Its centered one-hot vector is $g$. Using all numerically nonzero singular directions, set
\begin{equation}
 a=U\Sigma^{\dagger}V^{\mathsf T}g,\qquad u=N_ea.
 \label{eq:direction}
\end{equation}
This is the minimum-norm least-squares solution $a=(H^{\mathsf T})^\dagger g$; $u$ is not separately normalized. It fits class offsets using weights alone. We determine the largest resulting class offset, with index $k$, and require positive gaps
\begin{equation}
 \delta_{kj}=u^{\mathsf T}\widetilde P(w_k-w_j)>0
 \quad\text{for all }j\neq k.
\end{equation}
The construction is defined for directions satisfying this condition. All seven B/32 releases satisfy it at the fixed class draw; the condition is not automatic for an arbitrary classifier.

\subsection{Choosing a magnitude from a feature bound}
For each competitor $j$, define
\begin{equation}
 q_{kj}=\widetilde P(w_k-w_j),\quad
 C_d=I_d-\mathbf{1}_d\mathbf{1}_d^{\mathsf T}/d.
\end{equation}
The class-$k$ margin numerator of the protected public-base interpolation is
\begin{equation}
 z\operatorname{diag}(\widetilde\gamma)q_{kj}
 +\widetilde b^{\mathsf T}q_{kj}
 +\lambda\tau\delta_{kj}.
 \label{eq:margin}
\end{equation}
Because $z$ is centered, \cref{eq:lnbound} bounds its first term below by $-\sqrt d\|C_d\operatorname{diag}(\widetilde\gamma)q_{kj}\|_2$.

\begin{proposition}[Sufficient proxy condition]
Suppose all $\delta_{kj}$ are positive. If
\begin{equation}
\tau > \max_{j\neq k}
\frac{[\sqrt d\|C_d\operatorname{diag}(\widetilde\gamma)q_{kj}\|_2
-\widetilde b^{\mathsf T}q_{kj}]_+}{\lambda\delta_{kj}},
\label{eq:magnitude}
\end{equation}
then class $k$ has a strictly positive margin over every competitor for every normalized hidden state in the public-base interpolation model.
\end{proposition}
\begin{proof}
Substitution of the lower bound into \cref{eq:margin} makes each pairwise margin positive. The common feature normalization does not alter their signs.
\end{proof}

We set the magnitude to twice the bound and check the positive margins after FP32 materialization. This uses the same rule for every task, with no image-based fitting. The interpolation coefficient is fixed at $0.5$ for construction, including releases later evaluated in eight-expert merging.

\subsection{Analyzing a fixed merged model}
The same feature bound can be applied to the actual merged parameters after publication. For a candidate class $k$, let $q_{kj}=P_m(w_k-w_j)$ and compute
\begin{equation}
 L_{kj}=b_m^{\mathsf T}q_{kj}
 -\sqrt d\|C_d\operatorname{diag}(\gamma_m)q_{kj}\|_2.
 \label{eq:actualbound}
\end{equation}
If $\min_{j\neq k}L_{kj}>0$, the model predicts class $k$ for every normalized hidden state in exact arithmetic. The proof follows the same centered-feature bound as \cref{eq:margin}. This conclusion holds after arbitrary changes to preceding blocks, provided the final gain, bias, projection, and classifier stay fixed. It does not cover recovery that changes those final parameters. A negative lower bound is inconclusive.

We apply this test to the stored merged readouts after release. Partner weights enter this analysis, but not the construction. The bound concerns real-arithmetic evaluation of those stored tensors.
